% ATTEMPT_STATUS: unresolved
% TOP_PROBLEM: {"problemId": "problem.top500-omission-w045049-stationary-integral-varifold-singular-set", "canonicalTitle": "Measure-zero singular set of stationary integral varifolds", "releaseRank": 414, "primaryDomain": "analysis_pde", "primaryDomainLabel": "Analysis & PDE", "displayStatus": "open_with_solved_subcases", "releaseStatus": "open_with_solved_subcases"}
% !TeX program = lualatex
\documentclass[11pt]{article}
\usepackage[margin=25mm]{geometry}
\usepackage{fontspec}
\setmainfont{Latin Modern Roman}
\usepackage{amsmath,amssymb,mathtools}
\usepackage{unicode-math}
\setmathfont{Latin Modern Math}
\usepackage{xurl,hyperref}
\hypersetup{colorlinks=true,urlcolor=blue}
\setcounter{secnumdepth}{0}
\setlength{\parindent}{0pt}
\setlength{\parskip}{5pt}
\setlength{\emergencystretch}{3em}
\begin{document}
\section*{\texorpdfstring{Measure-zero singular set of stationary integral varifolds}{Measure-zero singular set of stationary integral varifolds}}\label{414}
\subsection{Definitions and mathematical statement}
An integral $m$-varifold in $U\subset{\mathbb{R}}^{m+n}$ is a Radon measure on position--tangent-plane pairs represented by a countably $m$-rectifiable set with locally integrable positive integer multiplicity. It is stationary when its first variation vanishes:
\[
 \int\operatorname{div}_P X(x){\,\mathrm{d}} V(x,P)=0
 \quad\text{for every }X\in C_c^1(U;{\mathbb{R}}^{m+n}).
\]
A regular point of $\operatorname{spt}V\cap U$ has a neighborhood on which $V$ is an integer multiple of a smooth embedded minimal $m$-submanifold. Let $\operatorname{Sing}(V)$ be the remaining support points. The conjecture is $\mathcal H^m(\operatorname{Sing}(V))=0$ for all positive $m,n$. Stationarity alone is assumed: area minimization, stability, and multiplicity one are not extra hypotheses. Rectifiability almost everywhere is weaker than having these open smooth neighborhoods.
\par\smallskip\textit{Formulation and concept references:} \href{https://arxiv.org/html/2503.00651}{[S1]}.

\subsection{Short English statement}
Does every stationary integral varifold have a singular set of zero measure in its own dimension, without any stability or minimizing assumption?
\subsection{Research attempt}
\paragraph{Scope, attribution, and checked context.}
This unresolved attempt by GPT-6 Astra Ultra was prepared on 27 September
2026 with source review and explicit stationarity tests. Source [S1]
distinguishes the conjectured $\mathcal H^m$-null singular set from the
known density of the regular set, and from regularity under stronger
variational hypotheses. Those distinctions are retained. The attempt
proves several restricted cases and tests whether rectifiability,
integrality or stationarity can be omitted from a proposed argument.
No area-minimizing or stability assumption is added.

\paragraph{What density one already gives.}
We use the classical Allard regularity theorem in the following standard
stationary integral setting: if a blow-up at a point is a
multiplicity-one plane, sufficiently small mass and tilt excess at
small scales give a smooth embedded minimal graph near that point.
The $C^{1,\alpha}$ conclusion is upgraded to smoothness by the minimal
surface equation. This is an established input recorded in [S1],
not an elementary theorem proved here.

At almost every point of an integral varifold, rectifiability gives
a tangent plane with an integer multiplicity $Q$. At points with
$Q=1$, that blow-up meets the multiplicity-one hypothesis.
Consequently the singular points of density one form a
$\|V\|$-null set. The unresolved part concerns points of
higher multiplicity; dividing the tangent plane by $Q$
does not divide the nearby varifold into $Q$ smooth sheets.

For stationary integral varifolds, monotonicity gives a lower
mass-density bound at every support point, by approximating
it with rectifiable points of integer density and applying
the monotonicity inequality. The usual covering argument
then makes $\mathcal H^m$ on the support controlled locally
by $\|V\|$. In particular replacing a mass-null exceptional
set by an $\mathcal H^m$-null one here does not discard
an additional large invisible part of the support.

\paragraph{A useful constant-multiplicity subcase.}
Suppose $p$ has tangent varifold $Q|P|$ and there is a neighborhood
$U_p$ in which the rectifiable multiplicity of $V$ equals
the same integer $Q$ almost everywhere. Then
$W=V/Q$ is integral and stationary on $U_p$ and has
multiplicity-one tangent $|P|$ at $p$. Allard regularity
applies to $W$ and proves that $V$ is $Q$ times a smooth
minimal graph near $p$. Thus $p$ is regular in precisely
the sense of the catalogue.

The hypothesis is local equality of the multiplicity,
not just the statement that its approximate limit at
$p$ is $Q$. A measurable integer-valued function is
approximately continuous almost everywhere, but need
not be constant on any open neighborhood there.
Small sets of other multiplicities may persist at every
scale. Replacing approximate constancy by open local
constancy would silently assume the main regularity issue.

\paragraph{Stationarity forbids an easy multiplicity counterexample.}
Let the support be one affine plane $P$, and write
$V=\theta\,\mathcal H^m\!\llcorner P$ with locally
integrable multiplicity $\theta$. Test stationarity
against vector fields tangent to $P$, compactly
supported in a connected plane domain. In coordinates,
\[
              \int_P\theta\,\partial_i\phi=0
              \quad\text{for all }\phi\in C_c^1(P)
              \text{ and all }i.
\]
Every distributional derivative of $\theta$ is zero.
Mollification makes this explicit: on smaller connected
domains, $\theta*\rho_\epsilon$ has zero ordinary
gradient and hence is constant; passage to the limit
shows that $\theta$ is almost everywhere constant.
If $V$ is integral, that constant is a positive integer.

Thus choosing multiplicity two on a positive-measure
Cantor subset of a plane and multiplicity one elsewhere
does not yield a stationary counterexample. Rectifiability
and integrality alone allow that construction, but its
tangential first variation is nonzero. This is a direct
test of a natural candidate that fails an actual hypothesis.

\paragraph{Finite unions of planes and the natural dimension barrier.}
A finite sum $V=\sum_{j=1}^N Q_j|P_j|$ of affine $m$-planes
with positive integer weights is stationary: the integral
of the tangential divergence on each whole plane vanishes.
Combine coincident planes first. Away from intersections
the varifold is a fixed integer multiple of a smooth plane.
The singular set is contained in the finite union of
$P_i\cap P_j$, $i\ne j$. Every nonempty such intersection
has dimension at most $m-1$, so its $\mathcal H^m$ measure
is zero.

This is a complete family with genuine crossing
singularities. Two transverse hyperplanes in
$\mathbb R^{m+1}$ meet in an $(m-1)$-plane, showing why
stationarity alone permits codimension-one singular
sets. It would be incorrect to import the much
stronger singular-set bounds of minimizing currents
into this problem.

Nor can one immediately obtain a positive-measure
singular set by taking infinitely many whole planes
with integer weights through a compact region.
If each such plane meets a fixed ball $B_r$,
its intersection with $B_{2r}$ contains an
$m$-ball of radius at least $r$. Each contributes
at least $\omega_mr^m$ mass. Infinitely many
distinct summands therefore violate local finite
mass. A more complicated construction would need
to evade this simple mass obstruction.

\paragraph{Dropping integrality produces a sharp failure test.}
Enumerate distinct rational heights $a_j$ dense in $(-1,1)$,
and let
\[
 V=\sum_{j\geq1}2^{-j}
             \bigl|\mathbb R^m\times\{a_j\}\bigr|
 \quad\text{in }\mathbb R^{m+1}.
\]
This is a rectifiable varifold with locally finite mass:
on a compact ball, each plane contributes at most the
volume of an $m$-ball of the same radius, and the
weights sum to one. Each summand is stationary, and
the absolute summability permits interchange with
the first-variation integral, so $V$ is stationary.

Its support is the entire closed slab
$\mathbb R^m\times[-1,1]$. No point in the open slab
has a neighborhood supported on one smooth embedded
$m$-submanifold. Thus the singular set contains an
open $(m+1)$-dimensional region, with nonzero
indeed locally infinite $\mathcal H^m$ measure.
The example fails the problem's integrality
hypothesis because the multiplicities are
$2^{-j}$. It is not a counterexample to the
catalogue assertion. It explains exactly why
the preceding integer mass obstruction matters.

The same construction in higher codimension is
obtained by adjoining zero coordinates. No
conclusion about the integral case follows from
this extension.

\paragraph{Why an open dense regular set is insufficient.}
The regular set is relatively open by definition.
Its density in the support is a known regularity
consequence discussed in [S1]. Neither topological
fact gives the required measure statement.
A closed nowhere-dense fat Cantor subset of an
interval has positive Lebesgue measure, while its
complement is open and dense. Products give the
same distinction in dimension $m$. These are
measure-theoretic tests, not examples of stationary
integral varifolds with such singular sets.

A successful argument must exploit more than
the ability to find one smooth neighborhood in
every ball. It would need a quantitative statement
showing that smooth neighborhoods exhaust almost
all of the mass, or a regularity mechanism at
almost every higher-multiplicity planar tangent.
The local constant-multiplicity proposition above
is one sufficient mechanism, but stationarity
has not been shown here to imply its hypothesis.

\paragraph{Verification, first gap, and outcome.}
The examples were checked directly in the first-variation
formula; none relies on an orientation or current
boundary. The infinite planar sum has finite mass
but deliberately noninteger weights. The integer
planar calculation accounts for coincident planes
and for their intersections. The distinction
between mass measure, support and Hausdorff
measure remains explicit throughout.

The first missing step is higher-multiplicity
regularity on a set of full $m$-dimensional
measure under stationarity alone. Tangent-plane
existence almost everywhere and integrality
do not by themselves supply open smooth
neighborhoods.
\textbf{UNRESOLVED.} The density-one and locally
constant-multiplicity cases follow from the
stated regularity input, and the elementary
countertests identify failures of weaker
hypotheses. No general null-singular-set theorem
or integral counterexample is claimed.

\subsection{Sources}
\begin{itemize}
\item[S1]Brena, Decio and De Lellis, Remarks and Conjectures on Stationary Varifolds (2025).
\url{https://arxiv.org/html/2503.00651}.
\textit{Formulation source.}
\end{itemize}
\end{document}
